% =====================================================================
%  Quatrième de couverture — page imposée par le modèle ISIMS.
%  Ordre du modèle : الخلاصة, puis Résumé, puis Abstract, puis les mots clés dans les trois langues.
%
%  Cette page exige XeLaTeX (voir preamble.tex) : l'arabe ne se compose pas sous pdfLaTeX.
% =====================================================================
\clearpage
\thispagestyle{empty}
% Comme la couverture, cette page est imposée et doit tenir sur une seule feuille : avec
% l'interligne 1,5 du corps, les trois résumés et le pied de page débordaient.
\singlespacing

\begin{center}
  {\bfseries\large Système intelligent de suivi de livraison}
\end{center}
\vspace{-0.2cm}
\noindent\textcolor{isimsblue}{\rule{\textwidth}{1pt}}

\begin{center}
  {\Large\bfseries\itshape Mohamed Aziz Hadjkacem}
\end{center}
\vspace{-0.2cm}
\noindent\textcolor{isimsblue}{\rule{\textwidth}{1pt}}

\vspace{0.3cm}
% Le guide de rédaction plafonne chaque résumé à 10 lignes (règle 6) : à ce format, en corps 12
% normal, les trois versions ci-dessous tiennent sans avoir besoin de réduire la police.

% --- Arabe -----------------------------------------------------------------
% Le titre est aligné à droite et le paragraphe composé de droite à gauche : dans le modèle, tout le
% bloc arabe est en RTL, y compris la ponctuation finale.
\begin{Arabic}
\begin{flushright}
\textbf{\underline{الخلاصة}} : تقدّم هذه المذكّرة منصّة لتتبّع عمليات التوصيل في الكيلومتر الأخير،
قابلة للاستعمال من قِبَل عدّة مؤسسات ومندمجة مع أنظمة تخطيط مواردها. تغطّي المنصّة استيراد الطلبات
من نظام تخطيط الموارد، وتخطيط الجولات، والتنفيذ الميداني عبر تطبيق جوّال يعمل دون اتصال بالإنترنت،
وإثبات التسليم، وإدارة الإرجاع. تعتمد البنية على الخدمات المصغّرة مع عزل البيانات حسب المخطّط،
ومصادقة مفوّضة عبر OAuth2/OIDC، ومُهايئ محايد لأنظمة تخطيط الموارد تم التحقّق منه على Odoo
وERPNext. تمّت النمذجة باستخدام UML وفق منهجية Scrum.
\end{flushright}
\end{Arabic}

\vspace{0.15cm}

% --- Français --------------------------------------------------------------
\noindent\textbf{\underline{Résumé}} : Ce mémoire présente une plateforme de suivi de livraison du
dernier kilomètre, exploitable par plusieurs entreprises clientes et intégrée à leurs progiciels de
gestion. Elle couvre l'import des commandes depuis l'ERP, la planification des tournées, l'exécution
sur le terrain via une application mobile fonctionnant hors connexion, la preuve de livraison et la
gestion des retours. L'architecture, orientée microservices, isole les données par schéma, délègue
l'authentification en OAuth2/OIDC et intègre un adaptateur ERP agnostique validé sur Odoo et
ERPNext. La modélisation a été réalisée en UML et le développement conduit selon la méthode Scrum.

\vspace{0.15cm}

% --- Anglais ---------------------------------------------------------------
\noindent\textbf{\underline{Abstract}}: This thesis presents a last-mile delivery tracking platform,
usable by several client companies and integrated with their enterprise resource planning systems.
It covers order import from the ERP, route planning, field execution through an offline-capable
mobile application, proof of delivery and returns management. The microservice-oriented architecture
isolates data by schema, delegates authentication to OAuth2/OIDC, and integrates an ERP-agnostic
adapter validated against both Odoo and ERPNext. Conceptual modelling was carried out using UML and
development followed the Scrum agile method.

\vspace{0.15cm}

\begin{Arabic}
\begin{flushright}
\textbf{\underline{الكلمات المفاتيح}} : التوصيل في الكيلومتر الأخير، التتبّع الآني، تعدّد
المستأجرين، تكامل أنظمة تخطيط الموارد، الخدمات المصغّرة، العمل دون اتصال، UML.
\end{flushright}
\end{Arabic}

\vspace{0.15cm}

\noindent\textbf{Mots clés} : livraison du dernier kilomètre, suivi en temps réel, multi-tenance,
intégration ERP, Odoo, ERPNext, microservices, mode hors ligne, UML.

\vspace{0.15cm}

\noindent\textbf{Keywords}: last-mile delivery, real-time tracking, multi-tenancy, ERP integration,
Odoo, ERPNext, microservices, offline mode, UML.

% \vfill poussait le pied de page tout en bas de la page disponible et ne lui laissait plus la
% place de s'y tenir : TeX le renvoyait alors, seul, sur une seconde feuille par ailleurs vide.
\vspace{0.5cm}

% --- Pied de page institutionnel (repris du modèle) ---
\noindent\textcolor{isimsblue}{\rule{\textwidth}{0.8pt}}\\[4pt]
\noindent
\begin{minipage}[c]{0.36\textwidth}
  \footnotesize
  www.isimsf.rnu.tn~\raisebox{-2pt}{\includegraphics[height=10pt]{ico-globe}}\\[4pt]
  % Coupure imposée : laissée au fil du texte, l'adresse cassait après « route de » et rejetait
  % l'icône seule sur une troisième ligne.
  Pôle technologique de Sfax\\
  route de Tunis km 10, B.P. 242 - 3021 Sfax~\raisebox{-2pt}{\includegraphics[height=10pt]{ico-pin}}
\end{minipage}
\hfill
\begin{minipage}[c]{0.20\textwidth}
  \centering
  \includegraphics[width=0.9\linewidth]{isims-logo-pied}
\end{minipage}
\hfill
\begin{minipage}[c]{0.30\textwidth}
  \raggedleft\footnotesize
  +216 74 862 432~\raisebox{-2pt}{\includegraphics[height=10pt]{ico-tel}}\\[4pt]
  +216 74 862 233~\raisebox{-2pt}{\includegraphics[height=10pt]{ico-fax}}
\end{minipage}
